\documentclass[a4paper]{scrartcl}
\usepackage[
    typ=kl,
    klausurtyp=klausur,
    fach=Informatik,
    lerngruppe=XX,
    erwartungshorizontAnzeigen,
    loesungen=folgend,
    module={Aufgabenpool}
]{schule}

\date{xx.xx.xxxx}
\title{Beispielaufgabe 1}

\begin{document}
    \begin{aufgabenpoolHinweis}
        Dieses ist nur eine Beispielaufgabe, die ohne weitere Optionen auskommt.
    \end{aufgabenpoolHinweis}

    \begin{aufgabe}[subtitle=Erste Beispielaufgabe]
        Blah blah blub.

        \begin{teilaufgaben}
            \teilaufgabe[2] Male ein X.
                \erwartung{malt ein X}{2}

            \teilaufgabe[3] Schreibe einen Satz.
                \erwartung{schreibt einen Satz}{2}
                \erwartung{achtet auf die Rechtschreibung}{1}
        \end{teilaufgaben}
    \end{aufgabe}
    \begin{loesung}
        \begin{teilaufgaben}
            \teilaufgabe Ein gemaltes X ist in \autoref{fig:x} zu sehen.
                \begin{figure}[ht]
                    \centering
                    \input{\getBasedir/zeichnung}
                    \caption{Ein x}
                    \label{fig:x}
                \end{figure}
            \teilaufgabe Eine möglicher Satz.
        \end{teilaufgaben}
    \end{loesung}
\end{document}
